% Step 3.1: FFR and robustness sweeps (validation split).
% Numbers from results/step3.1_ffr_val20_v60.json, results/step3.1_ffr_extraction_floor.json,
% results/step3.1_sweeps.json (all settings, Gate 1 method: 4DRGS initialization, calibrated radius, fixed geometry).
% Bib keys: thakur2026punch, young1973flow.

\paragraph{FFR.}
FFR is computed with the Young--Tsai model~\cite{young1973flow} from the
recovered geometry and flow, without ground truth: the lesion is located
at the largest relative radius deficit against a reference interpolated
from the neighboring segments, and its length is taken from the width of
the deficit. Because the bolus is imaged at rest, the hyperemic flow is the
recovered rest flow times the coronary flow reserve (CFR). With the CFR of
each case, FlowSplat reaches an FFR mean absolute error of 0.030 and 90\,\%
accuracy at the 0.80 threshold on 20 validation cases (60 views); with a
single population CFR of 2.75 the error rises to 0.079 (85\,\%), and the
same population CFR applied to the true geometry and flow already gives
0.071, so CFR uncertainty, not reconstruction, dominates in that setting.
An oracle decomposition shows that the FFR error stems from the lesion
geometry: with the true lesion and our flow the error is 0.012, with our
lesion and the true flow 0.034. The lesion extraction itself contributes
0.011 on the ground-truth profiles of all 200 cases.

\paragraph{Robustness.}
Table~\ref{tab:sweeps} varies one factor at a time around the 60-view
setting. The flow error stays between 1.7\,\% and 7.8\,\% across 30 to 133
views (with the radius calibration repeated per view count), dispersion up
to Péclet numbers of about 240, more than an order of magnitude below the
regime reported by PUNCH~\cite{thakur2026punch}, rest velocities from 0.10
to 0.40\,m/s, out-of-plane angles from 0 to 60\,deg, and four independent
reconstructions of the same data (per-case standard deviation 0.8 points).
The method is not robust to strong detector noise (flow error 20.9\,\%):
the reconstructed lumen becomes thinner and the calibration, obtained at
the default noise level, no longer compensates it. The FFR error follows
the lesion geometry; it doubles at 30 views (0.063) and fails when the
lesion finder picks a false narrowing, which happened in 3 of about 230
reconstructions.

\begin{table}[t]
\centering
\small
\caption{Robustness on the validation split (one factor varied, all others
at the 60-view setting in bold). Flow: mean relative error of $Q$;
diameter: mean stenosis diameter error; FFR: mean absolute error with the
case CFR. 20 validation cases for the views, 5 otherwise.
n.v.: ground-truth FFR not valid at that flow.}
\label{tab:sweeps}
\begin{tabular}{llccc}
\toprule
Factor & Level & Flow err. & Diam. err. (mm) & FFR MAE \\
\midrule
Views & 30 / \textbf{60} / 133 & 7.8 / \textbf{5.3} / 3.8\,\% & 0.17 / \textbf{0.10} / 0.11 & 0.063 / \textbf{0.030} / 0.040 \\
Noise & low / high & 6.1 / 20.9\,\% & 0.11 / 0.16 & 0.083 / -- \\
Dispersion $D$ & $10^{-5}$ / $10^{-4}$\,m$^2$/s & 3.6 / 4.5\,\% & 0.12 / 0.10 & 0.070 / 0.073 \\
Rest velocity & 0.10 / 0.25 / 0.40\,m/s & 3.3 / 4.4 / 6.3\,\% & 0.15 / 0.06 / 0.14 & 0.048 / n.v. / n.v. \\
Out of plane & 0 / 30 / 60\,deg & 4.5 / 1.7 / 4.9\,\% & 0.12 / 0.05 / 0.15 & 0.081 / 0.040 / 0.073$^\ast$ \\
Repeats & 4 runs & 3.3--4.5\,\% & 0.08--0.10 & 0.060--0.068 \\
\bottomrule
\end{tabular}

\vspace{2pt}
{\footnotesize $^\ast$ without one case with a misplaced lesion (0.531 with it). Noise high: FFR errors of geometry and flow cancel (0.025), not reported.}
\end{table}
